% Created 2025-03-06 Thu 12:44
% Intended LaTeX compiler: pdflatex
\documentclass[11pt]{article}
\usepackage[utf8]{inputenc}
\usepackage[T1]{fontenc}
\usepackage{graphicx}
\usepackage{longtable}
\usepackage{wrapfig}
\usepackage{rotating}
\usepackage[normalem]{ulem}
\usepackage{amsmath}
\usepackage{amssymb}
\usepackage{capt-of}
\usepackage{hyperref}
\date{\today}
\title{Managing my website with Org Mode: a Journey into Emacs}
\hypersetup{
 pdfauthor={},
 pdftitle={Managing my website with Org Mode: a Journey into Emacs},
 pdfkeywords={},
 pdfsubject={},
 pdfcreator={Emacs 30.1 (Org mode 9.7.22)}, 
 pdflang={English}}
\begin{document}

\maketitle
\tableofcontents

\section{Foreword}
\label{sec:orgd93711e}
\subsection{A Warning for Emacs Wizards}
\label{sec:orgba54347}
My skill in emacs and org mode is very much limited compared to most, so forgive me if I blaspheme -.-
\subsection{Dotfiles}
\label{sec:org1fb66cc}
You can find my dotfiles \href{https://kirawano.github.io/doom-kira}{here}, where they'll hopefully be updated with the changes mentioned in this blog by the time it's pushed to main.
\section{Management of Style with org-mode}
\label{sec:org4dd2793}
This journey started with me wanting to have a better experience making websites in a nicer-to-use text formatting style. While writing an essay for school, I noticed you could export .org files as html files, so the rusted gears in my head started turning in YARB-like fashion.

After a bit of searching, I added the following snippet to style.org.css in the root directory of this repo:

\begin{verbatim}
#+OPTIONS: org-html-head-include-default-style:nil


#+HTML_HEAD: <link rel="stylesheet" type="text/css" href="../style.css"> 
#+HTML_HEAD: <div class="topnav">
#+HTML_HEAD: <a class="home" href="https://kirawano.github.io">~</a>
#+HTML_HEAD: <a class="proj" href="https://kirawano.github.io/projects.html">Projects</a>
#+HTML_HEAD: <a class="YARB" href="https://kirawano.github.io/YARB.html">YARB</a>
#+HTML_HEAD: <a class="CONT" href="https://kirawano.github.io/contact.html">Contact</a>
#+HTML_HEAD: <a class="setp" href="https://kirawano.github.io/setup.html">setup</a>
#+HTML_HEAD: <a class="catd" href="https://kirawano.github.io/catdump.html">cat dump</a>
#+HTML_HEAD: <a class="x8831" href="https://kirawano.github.io/88x31.html">88x31</a>
#+HTML_HEAD:     </div>
#+HTML_HEAD:     <h1 id = 'titleText' style = 'color:rgb(204,0,255);'></h1>
#+HTML_HEAD:     </head>
\end{verbatim}

To fully utilize this, I created a tree strucutre of style.css.org files in each directory of the repo, each inheriting the org file of the super directory.
\end{document}
